\pagebreak
\section{Constriction Resistance}
Constriction resistance refers to the additional electrical resistance that arises at the interface between two conductive bodies due to the discrete nature of their mechanical contact. Although surfaces may appear smooth at the macroscopic scale, they are inherently rough when observed microscopically, consisting of asperities and voids, see \cref{fig:CR}. As a result, electrical contact is established only at a limited number of contact points rather than across the entire apparent contact area.
At these asperity contacts, the electrical current is forced to converge and pass through highly localized regions, leading to a constriction of the current flow. This geometric restriction increases the effective resistance compared to an ideal continuous conductor. The total constriction resistance is therefore governed by both the number of asperities and their individual contact sizes, which is strongly influenced by mechanical loading. 
As the load between the conductive surfaces increases, additional asperities come into contact and existing contact areas increase. This results in a larger effective region and a corresponding decrease in constriction resistance. Conversely, at low contact forces, only a few asperities contribute to conduction, leading to higher resistance.
This dependency enables constriction resistance to be exploited as a sensing mechanism, where mechanical stimuli are transduced into measurable changes in electrical resistance \cite{G.O.A.T}.

\begin{figure}[H]
\centering
\includestandalone[mode=buildmissing]{Standalone/Drawings/CR-tikz}
    \caption{Illustration of microscopic view between two surfaces creating contact with multiple asperities. Top figure with low force applied, and bottom with high force applied.}
    \label{fig:CR}
\end{figure}